%% ============================================================
%% shared/services/networking/hybrid_connectivity.tex
%% Hybrid Connectivity + VPC Peering, Transit Gateway, PrivateLink
%% Classification: BOTH (VPC Peering = SAA; TGW + PrivateLink deep = SAP)
%% ============================================================

\servicesection{Hybrid \& VPC Connectivity}{\bothtag}{%
  Connecting VPCs, on-premises networks, and services securely.}

% ─── TikZ: VPC Peering (non-transitive) vs Transit Gateway ────
\subsection{VPC Connectivity Options}

\begin{figure}[h]
\centering
\begin{tikzpicture}[
  vpc/.style={draw, circle, minimum size=1.2cm, font=\small\bfseries, align=center},
  peer/.style={draw, dashed, thick, color=saacolor},
  tgw/.style={draw, diamond, minimum width=1.4cm, minimum height=1.4cm,
               fill=awsorange!20, draw=awsorange, font=\small\bfseries},
  nope/.style={draw, very thick, color=warncolor, -}
]

%% ─── LEFT SIDE: VPC Peering (non-transitive) ───
\node[font=\bfseries, color=awsdark] at (0, 3.2) {VPC Peering (non-transitive)};

\node[vpc, fill=saacolor!15] (A) at (-1.5, 1.8) {A};
\node[vpc, fill=saacolor!15] (B) at ( 1.5, 1.8) {B};
\node[vpc, fill=saacolor!15] (C) at ( 0,   0.2) {C};

\draw[peer] (A) -- (B) node[midway, above, font=\tiny] {peered};
\draw[peer] (B) -- (C) node[midway, right, font=\tiny] {peered};

%% X mark showing A cannot reach C via B
\draw[nope, color=warncolor, dashed] (A) -- (C);
\node[color=warncolor, font=\bfseries] at (-0.5, 1.0) {\large\texttimes};
\node[font=\tiny, color=warncolor] at (-1.2, 0.6) {A $\not\to$ C};

%% ─── RIGHT SIDE: Transit Gateway (transitive) ───
\node[font=\bfseries, color=awsdark] at (7.5, 3.2) {Transit Gateway (hub-and-spoke)};

\node[tgw] (TGW) at (7.5, 1.5) {TGW};

\node[vpc, fill=sharedcolor!15] (D) at (5.5, 2.8) {A};
\node[vpc, fill=sharedcolor!15] (E) at (9.5, 2.8) {B};
\node[vpc, fill=sharedcolor!15] (F) at (7.5, 0.0) {C};
\node[vpc, fill=sapcolor!15]    (G) at (5.5, 0.0) {On-Prem};

\draw[->, thick, color=sharedcolor] (D) -- (TGW);
\draw[->, thick, color=sharedcolor] (E) -- (TGW);
\draw[->, thick, color=sharedcolor] (F) -- (TGW);
\draw[->, thick, color=sapcolor]    (G) -- (TGW) node[midway, left, font=\tiny] {VPN/DX};

\node[font=\tiny, color=sharedcolor] at (7.5, -0.8)
  {All VPCs can reach each other via TGW};

\end{tikzpicture}
\caption{Left: VPC Peering is not transitive (A cannot reach C via B). Right: Transit Gateway enables any-to-any routing including on-premises.}
\end{figure}

\begin{comparebox}[title={VPC Connectivity Options}]
\begin{tabular}{@{}p{3.5cm}p{5cm}p{3.5cm}@{}}
\toprule
\textbf{Method} & \textbf{Description} & \textbf{Key Constraint} \\
\midrule
VPC Peering          & Direct link between two VPCs; non-transitive         & No overlapping CIDRs; scales poorly at $>$20 VPCs \\
Transit Gateway (TGW) & Hub-and-spoke; transitive; cross-account; cross-region & Per-attachment + per-GB cost \\
AWS PrivateLink      & Expose a service to other VPCs without peering; traffic on AWS backbone & One-directional; consumer $\to$ provider \\
VPC Endpoints (Gateway) & Free private route to S3 and DynamoDB via route table & S3 and DynamoDB only \\
VPC Endpoints (Interface) & PrivateLink ENI in subnet; for most other AWS services & Per-AZ cost \\
\bottomrule
\end{tabular}
\end{comparebox}

\begin{gotcha}
\textbf{VPC Peering is NOT transitive.} If A peers with B and B peers with C, A cannot reach C via B. This is the most frequently tested VPC concept. Use \awsservice{Transit Gateway} when you need any-to-any connectivity across many VPCs or on-premises networks.
\end{gotcha}

% ─────────────────────────────────────────
\subsection{Hybrid Connectivity: Direct Connect \& VPN}

\begin{comparebox}[title={Hybrid connectivity options}]
\begin{tabular}{@{}p{3.5cm}p{5cm}p{3.2cm}@{}}
\toprule
\textbf{Service} & \textbf{Description} & \textbf{Speed / Note} \\
\midrule
AWS Direct Connect     & Dedicated private connection from on-premises to AWS; consistent low latency    & 1, 10, 100 Gbps \\
Site-to-Site VPN       & Encrypted IPsec tunnel over public internet; fast to provision                  & Up to 1.25 Gbps  \\
Client VPN             & Managed OpenVPN for remote user access                                         & Per-user \\
Direct Connect + VPN   & VPN over Direct Connect for encryption; combines privacy and encryption         & Production pattern \\
\bottomrule
\end{tabular}
\end{comparebox}

\begin{examtip}
\textbf{Resilient hybrid connectivity pattern:} use \textbf{Direct Connect as primary} (consistent, low-latency) and \textbf{VPN as failover}. Direct Connect alone is not redundant --- provision two DX connections from different providers or pair with VPN. For maximum resilience: two DX connections in different locations.
\end{examtip}

\begin{saadepth}
For SAA:
\begin{itemize}
  \item Direct Connect takes \textbf{weeks} to provision; VPN is \textbf{minutes}.
  \item Use VPC \textbf{Gateway Endpoints} for S3/DynamoDB access (free, no data processing charge vs NAT Gateway).
  \item Use VPC \textbf{Interface Endpoints} (PrivateLink) for other AWS services to avoid traffic leaving the VPC.
\end{itemize}
\end{saadepth}

\begin{sapdepth}
SAP networking depth:
\begin{itemize}
  \item \textbf{Transit Gateway + Direct Connect Gateway} --- connect a single TGW to multiple regions and on-premises over DX; use transit VIF (Virtual Interface).
  \item \textbf{Shared VPC via AWS RAM} --- share subnets from a central network account to spoke accounts; centralised routing without peering mesh.
  \item \textbf{Centralised egress via TGW + NAT Gateway} --- single NAT GW in a dedicated egress VPC; all spoke VPCs route through TGW to the NAT GW.
  \item \textbf{AWS Network Firewall} --- deploy in centralised inspection VPC; route all ingress/egress through it via TGW and Gateway Load Balancer.
  \item \textbf{Direct Connect SLA:} 99.9\% (1 connection) or 99.99\% (2 connections from different locations). Design for the SLA you need.
\end{itemize}
\end{sapdepth}

% ─────────────────────────────────────────
\subsection{Hybrid DNS with Route 53 Resolver}

\begin{keyfacts}
\begin{itemize}
  \item \textbf{Inbound Resolver Endpoint} --- allows on-premises DNS servers to resolve AWS private hosted zone records (on-prem $\to$ AWS DNS)
  \item \textbf{Outbound Resolver Endpoint} --- allows VPC instances to forward DNS queries to on-premises DNS (AWS $\to$ on-prem DNS)
  \item \textbf{Forwarding Rules} --- define which domains are forwarded to on-premises resolvers
  \item Deploy at least two endpoints across AZs for high availability
\end{itemize}
\end{keyfacts}

\begin{examtip}
On the SAP exam: hybrid DNS questions typically involve a company that has on-premises DNS (e.g., Active Directory) that needs to resolve AWS private hostnames (VPC DNS) and vice versa. The answer is always \textbf{Route 53 Inbound + Outbound Resolver Endpoints} with forwarding rules --- not a custom DNS on EC2.
\end{examtip}
